\begin{exercise}[Gaz de van der Waals]
On donne l’équation d’état d’un gaz réel de van der Waals, contenant 7 mol de gaz :

\begin{align*}
\left( P+ a \frac{n^2}{V} \right) \cdot (V-nb) = nRT
\end{align*}

où $a$ et $b$ sont des constantes caractéristique du fluide, et $R$, la constante des gaz parfaits. 

L'expression différentielle de l’énergie interne du gaz s’écrit :
\begin{align*}
d U = n C_{V,m} dT + a \frac{n^2}{V^2} dV
\end{align*}

\begin{questions}
\question Montrer qu’une des identités thermodynamiques permet d’écrire l’expression différentielle $dS(T,V)$.

\question $C_{V,m}$ est constant. En déduire, à une constante près, les fonctions d’état $U(T,V)$ et $S(T,V)$.
\end{questions}

\end{exercise}